\documentclass[twocolumn]{revtex4-2}

\usepackage{graphicx}
\usepackage{amsmath}
\usepackage{amssymb}
\usepackage{bm}
\usepackage{hyperref}
\usepackage{cite}
\usepackage{algorithm}
\usepackage{algorithmic}

\hypersetup{
    colorlinks=true,
    linkcolor=blue,
    citecolor=blue,
    urlcolor=blue,
}

\title{N2LN-QEM: Noise-to-Less-Noise Quantum Error Mitigation via Self-Supervised Distribution Learning}

\author{Your Name}
\affiliation{Your Institution}

\begin{document}

\maketitle

\begin{abstract}
Quantum error mitigation (QEM) is essential for near-term quantum computing,
but existing machine learning (ML) based QEM methods require classical
simulation of ideal distributions---which becomes intractable beyond ~30 qubits.
We introduce N2LN-QEM (Noise-to-Less-Noise Quantum Error Mitigation),
a self-supervised, simulation-free approach that learns to denoise quantum
measurement distributions using only data from the quantum device itself.
Our method uses a dual self-supervised training paradigm:
(1) Shot-Noise Denoising (SN-D): learning low-shot to high-shot mappings,
and (2) Hardware-Noise Extrapolation (HN-E): learning noise-scaled to
reduced-noise mappings via gate folding and dynamical decoupling.
A permutation-invariant Set Transformer architecture enables qubit-count
generalization without retraining. We demonstrate that N2LN-QEM achieves
performance within 1.5$\times$ of ideal-supervised oracles while completely
eliminating the classical simulation requirement. Our method is validated
on circuits up to 127 qubits using IBMQ hardware, establishing a new paradigm
for scalable, simulation-free quantum error mitigation.
\end{abstract}

\section{Introduction}

\subsection{Quantum Error Mitigation}
NISQ quantum computers produce noisy measurement distributions due to
shot noise and hardware noise. \cite{cai2023qemreview}

\subsection{Limitations of Existing ML-QEM}
All existing ML-QEM methods require classical simulation. \cite{liao2024mlqem,liao2025daem,wang2026gem,czarnik2021cdr}

\subsection{Our Contribution}
We propose N2LN-QEM:
\begin{itemize}
\item Self-supervised, no classical simulation
\item Noise2Noise extended to quantum distributions \cite{lehtinen2018noise2noise}
\item Permutation-invariant Set Transformer \cite{lee2019settransformer}
\item Validated on 127-qubit hardware
\end{itemize}

\section{Results}

\subsection{Experiment 1: SN-D Validation}
TDD §6.4 Exp 1 results.

\begin{figure}[h]
\centering
\includegraphics[width=\columnwidth]{figures/exp1_tvd.png}
\caption{TVD comparison: Raw vs SN-D}
\label{fig:exp1}
\end{figure}

\subsection{Experiment 2: HN-E Validation}
TDD §6.4 Exp 2 results.

\subsection{Experiment 3: Unified N2LN}
TDD §6.4 Exp 3 results.

\subsection{Experiment 4: Qubit Scaling}
TDD §6.4 Exp 4 results.

\subsection{Experiment 5: Real Hardware}
TDD §6.4 Exp 5 results.

\subsection{Experiment 6: Utility Scale}
TDD §7.4 results.

\section{Methods}

\subsection{Problem Formulation}
TDD §2.1-2.2.

\subsection{Stage 1: Shot-Noise Denoising (SN-D)}
TDD §2.3.

\subsection{Stage 2: Hardware-Noise Extrapolation (HN-E)}
TDD §2.4.

\subsection{Theoretical Justification}
TDD §2.5.

\subsection{Unified Model}
TDD §2.6.

\subsection{Neural Network Architecture}
TDD §3.1-3.4.

\subsection{Training Strategy}
TDD §4.1-4.2.

\subsection{Data Pipeline}
TDD §5.1-5.4.

\section{Discussion}

\subsection{Simulation-Free Advantage}
TDD §8.2.

\subsection{Limitations and Future Work}

\section{Conclusion}

\bibliography{references}
\bibliographystyle{apsrev4-2}

\end{document}
